\section{结论}
\label{sec:conclusion}

本文深入介绍了 \sysname 的设计、实现和部署。它是一个生产级系统，用于高效训练 MoE 模型。\sysname 利用通信高效方法，包括通信量更低的并行策略、算子间和算子内通信-计算重叠，以及结合调整后通信模式的通信压缩，从而释放高性能 GPU 的计算能力。在 1,440 张 NVIDIA Hopper GPU 上训练 352B MoE 模型时，\sysname 达到 1.41M tokens/s 吞吐，相比 Megatron-LM 提升 1.88$\times$。通过分享我们关于加速大规模 MoE 训练的见解，我们希望本工作能够启发未来研究。

% We emphasize the need for fault tolerance throughout the training process and implement a tailored robust training framework to locate and fix faults automatically. We provide a comprehensive set of monitoring tools for deep observability into system components and events, facilitating root cause identification for intricate anomalies. We believe that our work not only offers practical insights for those working on LLM training, but also paves the way for future research in this rapidly evolving field.
